\scanpage{153}
$X$. En d'autres termes, on suppose donné \deleted{un hom d'anneaux topologiques}\inserted{ des \deleted{\unsure{hom}} éléments}
\[
\pi_i\in\operatorname{End}\operatorname{id}_{\mathcal V},\quad\text{avec}\quad
\pi_i^2=\pi_i,\quad\pi_i\pi_j=\pi_j\pi_i=0\quad\text{si }i\ne j,
\]
\[
\deleted{k[[\pi_1,\ldots,\pi_n,\ldots]]/(\pi_i^2-\pi_i,\ \pi_i\pi_j,\ \deleted{\unsure{\pi_j}}\ \sum\pi_i-1}
\]
et $\sum\pi_i=1$. \margin{somme infinie !}
Alors les $\pi_i$ s'étendent à $\mathcal M$, qui devient une catégorie graduée, \deleted{\unsure{par}} ``somme'' de catégories $\mathcal M^{i}_{\deleted{\unsure{i}}}$, dont les objets sont appelés ``motifs\inserted{ effectifs} de poids $i$''.

On suppose de plus que dans $X\boxtimes Y$, on a
\[
\pi_n(X\boxtimes Y)=\sum_{i+j=n}\pi_i(X)\boxtimes\pi_j(Y);
\]
celà implique alors que
\diagram{\begin{tikzcd}
\mathcal M^i\times\mathcal M^j \arrow[r,dash] & \mathcal M^{i+j}
\end{tikzcd}}
par l'accouplement $\otimes$, et fournit donc un lien étroit entre les $\mathcal M^i$. Ainsi, le groupe\inserted{ de classes} $K$ de $\mathcal M$ est un \uline{$\lambda$-anneau gradué}.
\margin{L'objet unité $\Lambda(0)$ est de degré $0$. \uline{On va supposer que c'est le seul objet simple de ce degré}, et \uline{que \unsure{son end. soit} \inserted{\unsure{identifié à}} $\mathbf Q$}.}

\deleted{\unsure{On}} Je suppose maintenant fixé, de plus, un objet \deleted{\unsure{Z(1)}} $\Lambda(-1)$ de $\mathcal M$,\inserted{ de degré $\unsure{\rho|\beta}\geq1$,} tel que, posant \deleted{\unsure{\Lambda}} $M(1)=M\otimes\Lambda(1)$, \deleted{\unsure{le foncteur}} le
\mspage{152}
